% GENERATED FILE. Edit canonical lessons, then run npm run generate:books.
% canonical-source-hash: fnv1a64:39131793403833d2
% canonical-lessons: TE-C372-kathakudu, TE-C372-kathamsam, TE-C372-patra, TE-C372-kalpana, TE-C372-kattukatha

\chapter{How a Story Works}
\label{ch:te-how-a-story-works}
\hlchaptermodality{372}

\begin{chapteropening}
\textbf{By the end of this chapter:} \emph{I can say a storyteller, a plot, a character, imagination and a made-up story.}

Complete the last lesson of chapter 372: I can say a storyteller, a plot, a character, imagination and a made-up story.
\end{chapteropening}

\section[\texorpdfstring{\te{కథకుడు} (kathakuḍu)}{కథకుడు (kathakuḍu)}]{\te{కథకుడు} (kathakuḍu) — a storyteller; a narrator}

\label{lesson:TE-C372-kathakudu}



\begin{quote}
Before the new one: say the Telugu for a boon, then the Telugu for a curse.
\end{quote}

\subsection*{You'll want to know: \te{కథకుడు}}
\textbf{\te{కథకుడు}} — \emph{kathakuḍu} — \textquotedblleft{}a storyteller; a narrator\textquotedblright{}.

\begin{quote}
\te{కథకుడు} \te{పిల్లలకు} \te{ఒక} \te{వింత} \te{కథ} \te{చెప్పాడు}.
\end{quote}

\begin{quote}
\emph{kathakuḍu pillalaku oka vinta katha ceppāḍu.}
\end{quote}

\begin{quote}
The storyteller told the children a strange story.
\end{quote}

\te{కథ}, a story, with -\te{కుడు}: one who tells it.

\begin{cousinweb}
\textbf{\te{కథకుడు}} (\emph{kathakuḍu}) is borrowed from Sanskrit \emph{kathaka}, \textquotedblleft{}a narrator, one who recites a story\textquotedblright{} (Monier-Williams), built on \emph{kath}, \textquotedblleft{}to tell\textquotedblright{}.
\end{cousinweb}

\begin{grammarlens}[title={What you've built}]
One more word for telling what happened.
\end{grammarlens}

\subsection*{Your turn}
\begin{itemize}
\raggedright
  \item \emph{Say it:} \emph{kathakuḍu}
  \item \emph{Say it:} \emph{kathakuḍu}, once more
  \item \emph{Say it:} \emph{kathakuḍu pillalaku oka vinta katha ceppāḍu.} — The storyteller told the children a strange story.
  \item \emph{Recall:} say the Telugu for a boon, then the Telugu for a curse, then say \emph{kathakuḍu} again
\end{itemize}

\begin{tcolorbox}[breakable,colback=teal!4,colframe=teal!35!black,title={Before you move on}]
What does \te{కథకుడు} mean? (\textquotedblleft{}A storyteller; a narrator\textquotedblright{}.) Say it once more.
\end{tcolorbox}
\section[\texorpdfstring{\te{కథాంశం} (kathāṁśaṁ)}{కథాంశం (kathāṁśaṁ)}]{\te{కథాంశం} (kathāṁśaṁ) — a plot (of a story)}

\label{lesson:TE-C372-kathamsam}



\begin{quote}
Before the new one: say the Telugu for a curse, then the Telugu for a storyteller.
\end{quote}

\subsection*{You'll want to know: \te{కథాంశం}}
\textbf{\te{కథాంశం}} — \emph{kathāṁśaṁ} — \textquotedblleft{}a plot (of a story)\textquotedblright{}.

\begin{quote}
\te{ఆ} \te{సినిమా} \te{కథాంశం} \te{చాలా} \te{సులభం}.
\end{quote}

\begin{quote}
\emph{ā sinimā kathāṁśaṁ cālā sulabhaṁ.}
\end{quote}

\begin{quote}
That film's plot is very simple.
\end{quote}

\te{కథ}, a story, with a Sanskrit word for a part or a point.

\begin{cousinweb}
\textbf{\te{కథాంశం}} (\emph{kathāṁśaṁ}) is borrowed from Sanskrit: \emph{kathā}, \textquotedblleft{}a story\textquotedblright{}, and \emph{aṁśa}, \textquotedblleft{}a share, a part\textquotedblright{} (CDIAL 2).
\end{cousinweb}

\begin{grammarlens}[title={What you've built}]
One more word for telling what happened.
\end{grammarlens}

\subsection*{Your turn}
\begin{itemize}
\raggedright
  \item \emph{Say it:} \emph{kathāṁśaṁ}
  \item \emph{Say it:} \emph{kathāṁśaṁ}, once more
  \item \emph{Say it:} \emph{ā sinimā kathāṁśaṁ cālā sulabhaṁ.} — That film's plot is very simple.
  \item \emph{Recall:} say the Telugu for a curse, then the Telugu for a storyteller, then say \emph{kathāṁśaṁ} again
\end{itemize}

\begin{tcolorbox}[breakable,colback=teal!4,colframe=teal!35!black,title={Before you move on}]
What does \te{కథాంశం} mean? (\textquotedblleft{}A plot (of a story)\textquotedblright{}.) Say it once more.
\end{tcolorbox}
\section[\texorpdfstring{\te{పాత్ర} (pātra)}{పాత్ర (pātra)}]{\te{పాత్ర} (pātra) — a character (in a story); a role}

\label{lesson:TE-C372-patra}



\begin{quote}
Before the new one: say the Telugu for a storyteller, then the Telugu for a plot.
\end{quote}

\subsection*{You'll want to know: \te{పాత్ర}}
\textbf{\te{పాత్ర}} — \emph{pātra} — \textquotedblleft{}a character (in a story); a role\textquotedblright{}.

\begin{quote}
\te{ఆ} \te{నాటకంలో} \te{నాది} \te{చిన్న} \te{పాత్ర}.
\end{quote}

\begin{quote}
\emph{ā nāṭakaṁlō nādi cinna pātra.}
\end{quote}

\begin{quote}
I have a small part in that play.
\end{quote}

Plural: \te{పాత్రలు}, \emph{pātralu}.

Also a vessel or a pot: the Sanskrit word meant a vessel first.

\begin{cousinweb}
\textbf{\te{పాత్ర}} (\emph{pātra}) is borrowed from Sanskrit \emph{pātra}, \textquotedblleft{}a drinking vessel, a dish\textquotedblright{} (CDIAL 8055); in Telugu it is also a part in a play.
\end{cousinweb}

\begin{grammarlens}[title={What you've built}]
One more word for telling what happened.
\end{grammarlens}

\subsection*{Your turn}
\begin{itemize}
\raggedright
  \item \emph{Say it:} \emph{pātra}
  \item \emph{Say it:} \emph{pātra}, once more
  \item \emph{Say it:} \emph{ā nāṭakaṁlō nādi cinna pātra.} — I have a small part in that play.
  \item \emph{Recall:} say the Telugu for a storyteller, then the Telugu for a plot, then say \emph{pātra} again
\end{itemize}

\begin{tcolorbox}[breakable,colback=teal!4,colframe=teal!35!black,title={Before you move on}]
What does \te{పాత్ర} mean? (\textquotedblleft{}A character (in a story); a role\textquotedblright{}.) Say it once more.
\end{tcolorbox}
\section[\texorpdfstring{\te{కల్పన} (kalpana)}{కల్పన (kalpana)}]{\te{కల్పన} (kalpana) — imagination; fiction}

\label{lesson:TE-C372-kalpana}



\begin{quote}
Before the new one: say the Telugu for a plot, then the Telugu for a character.
\end{quote}

\subsection*{You'll want to know: \te{కల్పన}}
\textbf{\te{కల్పన}} — \emph{kalpana} — \textquotedblleft{}imagination; fiction\textquotedblright{}.

\begin{quote}
\te{ఆ} \te{కథ} \te{పూర్తిగా} \te{కల్పన}.
\end{quote}

\begin{quote}
\emph{ā katha pūrtigā kalpana.}
\end{quote}

\begin{quote}
That story is pure imagination.
\end{quote}

\begin{cousinweb}
\textbf{\te{కల్పన}} (\emph{kalpana}) is borrowed from Sanskrit \emph{kalpanā}, \textquotedblleft{}making, preparing\textquotedblright{} (Monier-Williams); in Telugu it is what the mind makes up.
\end{cousinweb}

\begin{grammarlens}[title={What you've built}]
One more word for telling what happened.
\end{grammarlens}

\subsection*{Your turn}
\begin{itemize}
\raggedright
  \item \emph{Say it:} \emph{kalpana}
  \item \emph{Say it:} \emph{kalpana}, once more
  \item \emph{Say it:} \emph{ā katha pūrtigā kalpana.} — That story is pure imagination.
  \item \emph{Recall:} say the Telugu for a plot, then the Telugu for a character, then say \emph{kalpana} again
\end{itemize}

\begin{tcolorbox}[breakable,colback=teal!4,colframe=teal!35!black,title={Before you move on}]
What does \te{కల్పన} mean? (\textquotedblleft{}Imagination; fiction\textquotedblright{}.) Say it once more.
\end{tcolorbox}
\section[\texorpdfstring{\te{కట్టుకథ} (kaṭṭukatha)}{కట్టుకథ (kaṭṭukatha)}]{\te{కట్టుకథ} (kaṭṭukatha) — a made-up story}

\label{lesson:TE-C372-kattukatha}



\begin{quote}
Before the new one: say the Telugu for a character, then the Telugu for imagination.
\end{quote}

\subsection*{You'll want to know: \te{కట్టుకథ}}
\textbf{\te{కట్టుకథ}} — \emph{kaṭṭukatha} — \textquotedblleft{}a made-up story\textquotedblright{}.

\begin{quote}
\te{అతను} \te{చెప్పింది} \te{అంతా} \te{కట్టుకథ}.
\end{quote}

\begin{quote}
\emph{atanu ceppindi antā kaṭṭukatha.}
\end{quote}

\begin{quote}
All he said was a made-up story.
\end{quote}

\te{కట్టు}, to tie or put together, and \te{కథ}: a story stitched together.

\begin{cousinweb}
\textbf{\te{కట్టుకథ}} (\emph{kaṭṭukatha}) is a Telugu compound: native \emph{kaṭṭu}, \textquotedblleft{}to tie, build\textquotedblright{} (DEDR 1147), and \emph{katha}, \textquotedblleft{}a story\textquotedblright{}, from Sanskrit.
\end{cousinweb}

\begin{grammarlens}[title={What you've built}]
One more word for telling what happened.
\end{grammarlens}

\subsection*{Your turn}
\begin{itemize}
\raggedright
  \item \emph{Say it:} \emph{kaṭṭukatha}
  \item \emph{Say it:} \emph{kaṭṭukatha}, once more
  \item \emph{Say it:} \emph{atanu ceppindi antā kaṭṭukatha.} — All he said was a made-up story.
  \item \emph{Recall:} say the Telugu for a character, then the Telugu for imagination, then say \emph{kaṭṭukatha} again
\end{itemize}

\begin{tcolorbox}[breakable,colback=teal!4,colframe=teal!35!black,title={Before you move on}]
What does \te{కట్టుకథ} mean? (\textquotedblleft{}A made-up story\textquotedblright{}.) Say it once more.
\end{tcolorbox}
